% Einheit 5 von 5 – Vergleichen, Ordnen, Runden (bank/brueche-dezimalzahlen/e5.jsonl, zusammenbau v0.1)
%% TODO Zweigzeile Teil 1: was hier gelernt wird (Satz in Schülersprache aus dem Einheitstitel) – Einheit 5
\einheitenkopf[e5]{Einheit 5 von 5 · Vergleichen, Ordnen, Runden}
\zweigzeile{ab Kl. 5, je nach Buch bis Kl. 6 · P10 oft}
\setcounter{aufgabe}{24}

%% TODO Titel: Ich-kann-Satz fehlt in der Bank (Platzhalter: Kettenname) – Nr. 25
%% TODO Anweisung: ein Satz über der ersten Teilaufgabe fehlt in der Bank – Nr. 25
\begin{aufgabe}{Welche Form?}
\begin{teile}
\teil Welche Form hat $\frac{7}{4}$? Kreuze an.\\ \kreuz{Bruch}\\ \kreuz{Dezimalzahl}\\ \kreuz{Prozent}\\ \kreuz{Potenz}
\end{teile}
\end{aufgabe}

%% TODO Titel: Ich-kann-Satz fehlt in der Bank (Platzhalter: Kettenname) – Nr. 26
%% TODO Anweisung: ein Satz über der ersten Teilaufgabe fehlt in der Bank – Nr. 26
\begin{aufgabe}{Vergleichen und Ordnen}
\begin{teile}
\teil Hänge Nullen an, bis beide gleich viele Stellen haben: $0{,}7$ und $0{,}42$. \leerfeld und \leerfeld
\end{teile}
%% TODO Darstellung für schwach fehlt in der Bank (brueche-dezimalzahlen-e5-k2-s1-v1; Abschnitt „Für schwache Schüler“)
\swz{Welche Zahl ist größer: $0{,}47$ oder $0{,}52$?}{}
%% TODO Darstellung für schwach fehlt in der Bank (brueche-dezimalzahlen-e5-k2-s1-v2; Abschnitt „Für schwache Schüler“)
\swz{Welche Zahl ist größer: $3{,}81$ oder $3{,}18$?}{}
%% TODO Darstellung für schwach fehlt in der Bank (brueche-dezimalzahlen-e5-k2-s1-v3; Abschnitt „Für schwache Schüler“)
\swz{Welche Zahl ist größer: $0{,}4$ oder $0{,}9$?}{}
%% TODO Darstellung für schwach fehlt in der Bank (brueche-dezimalzahlen-e5-k2-s1-v4; Abschnitt „Für schwache Schüler“)
\swz{Welche Zahl ist größer: $12{,}35$ oder $12{,}53$?}{}
%% TODO Darstellung für schwach fehlt in der Bank (brueche-dezimalzahlen-e5-k2-s1-v5; Abschnitt „Für schwache Schüler“)
\swz{Welche Zahl ist größer: $0{,}750$ oder $0{,}705$?}{}
\swfrage{Was bleibt gleich, was ändert sich?}
%% TODO Darstellung für schwach fehlt in der Bank (brueche-dezimalzahlen-e5-k2-s2-v1; Abschnitt „Für schwache Schüler“)
\swz{Welche Zahl ist größer: $0{,}8$ oder $0{,}79$?}{}
%% TODO Darstellung für schwach fehlt in der Bank (brueche-dezimalzahlen-e5-k2-s3-v1; Abschnitt „Für schwache Schüler“)
\swz{Ordne von klein nach groß: $0{,}4$; $0{,}38$; $0{,}405$; $0{,}39$.}{}
%% TODO Darstellung für schwach fehlt in der Bank (brueche-dezimalzahlen-e5-k2-s4-v1; Abschnitt „Für schwache Schüler“)
\swz{Runde $3{,}864$ auf Zehntel.}{}
%% TODO Darstellung für schwach fehlt in der Bank (brueche-dezimalzahlen-e5-k2-s5-v1; Abschnitt „Für schwache Schüler“)
\swz{Welche Zahl liegt genau in der Mitte zwischen $1{,}4$ und $1{,}5$?}{}
%% TODO Darstellung für schwach fehlt in der Bank (brueche-dezimalzahlen-e5-k2-s6-v1; Abschnitt „Für schwache Schüler“)
\swz{Setze $<$, $=$ oder $>$ ein: $\frac{3}{5}$ __ $0{,}65$}{}
%% TODO Darstellung für schwach fehlt in der Bank (brueche-dezimalzahlen-e5-k2-s7-v1; Abschnitt „Für schwache Schüler“)
\swz{Setze $<$, $=$ oder $>$ ein: $0{,}4$ __ $38\,\%$}{}
%% TODO Darstellung für schwach fehlt in der Bank (brueche-dezimalzahlen-e5-k2-s8-v1; Abschnitt „Für schwache Schüler“)
\swz{Setze $<$, $=$ oder $>$ ein: $0{,}5^2$ __ $0{,}5$}{}
\swz[3]{Welcher Wert ist der kleinste: $5{,}5$; $0{,}55$; $0{,}5^2$; $55\,\%$? \hfill (P10 2023 OS)}{}
\end{aufgabe}

%% TODO Titel: Ich-kann-Satz fehlt in der Bank (Platzhalter: Kettenname) – Nr. 27
%% TODO Anweisung: ein Satz über der ersten Teilaufgabe fehlt in der Bank – Nr. 27
\begin{aufgabe}{Zahl zwischen zwei Zahlen angeben (auch zwischen zwei Brüchen über Dezimalzahlen)}
%% TODO Darstellung für schwach fehlt in der Bank (brueche-dezimalzahlen-e5-k3-s1-v1; Abschnitt „Für schwache Schüler“)
\swz{Gib eine Zahl zwischen $0{,}7$ und $0{,}8$ an.}{}
\end{aufgabe}

%% TODO Titel: Ich-kann-Satz fehlt in der Bank (Platzhalter: Kettenname) – Nr. 28
%% TODO Anweisung: ein Satz über der ersten Teilaufgabe fehlt in der Bank – Nr. 28
\begin{aufgabe}{gemischte Liste ordnen (Bruch, Dezimalzahl, Prozent, Potenz einer Dezimalzahl)}
%% TODO Darstellung für schwach fehlt in der Bank (brueche-dezimalzahlen-e5-k4-s1-v1; Abschnitt „Für schwache Schüler“)
\swz{Ordne von klein nach groß: $\frac{1}{2}$; $0{,}45$; $48\,\%$; $0{,}7^2$.}{}
\end{aufgabe}

%% TODO Titel: Ich-kann-Satz fehlt in der Bank (Platzhalter: Kettenname) – Nr. 29
%% TODO Anweisung: ein Satz über der ersten Teilaufgabe fehlt in der Bank – Nr. 29
\begin{aufgabe}{Aussagen prüfen und die wahre ankreuzen}
\begin{teile}
\teil Welche Aussage ist wahr? Kreuze an.\\ \kreuz{$0{,}3 > 0{,}29$}\\ \kreuz{$\frac{1}{3} = 0{,}3$}\\ \kreuz{$3\,\% = 0{,}3$}
\end{teile}
\end{aufgabe}

%% TODO Titel: Ich-kann-Satz fehlt in der Bank (Platzhalter: Kettenname) – Nr. 30
%% TODO Anweisung: ein Satz über der ersten Teilaufgabe fehlt in der Bank – Nr. 30
\begin{aufgabe}{Vergleichen und Ordnen – Fehler finden, Begründen, Darstellungswechsel, Anwendung}
\swz[3]{Lea sagt: $0{,}28 > 0{,}3$, weil $28$ größer ist als $3$. Finde den Fehler.}{}
\swfrage{Warum ist $0{,}7$ größer als $0{,}69$? Begründe.}
\swz{Schreibe $0{,}65$ als gekürzten Bruch und in Prozent.}{}
\swz[3]{Beim 100-m-Lauf brauchen Ali $13{,}4$ s, Ben $13{,}38$ s und Can $13{,}45$ s. Wer ist am schnellsten, wer wird Dritter?}{}
\end{aufgabe}

\uebersichtskasten{Dezimalzahlen vergleichen: Stelle für Stelle von links – vorher Nullen anhängen, bis alle gleich viele Stellen haben. \quad 0,6 > 0,58 (0,60 > 0,58) \\ Runden: Die Stelle dahinter entscheidet – 0 bis 4 ab, 5 bis 9 auf. \quad 2,749 ≈ 2,7 (Zehntel), ≈ 2,75 (Hundertstel) \\ Mitte zweier Zahlen: beide addieren, durch 2 – oder am Zahlenstrahl eine Stelle feiner einteilen. \quad Mitte von 0,3 und 0,4: 0,35 \\ Verschiedene Darstellungen: alle in Dezimalzahlen umwandeln, dann ordnen. \quad 3/8 = 0,375 \quad 35 \% = 0,35 \quad 0,6² = 0,36 → 0,35 < 0,36 < 0,375}
